% =============================================
% File: chapters/methodology.tex
% Chương 3: Phương pháp thực hiện
% =============================================

\chapter{PHƯƠNG PHÁP THỰC HIỆN}
\label{chap:phuongphap}

Chương này trình bày toàn bộ quy trình và phương pháp thực hiện đề tài theo hai
nhánh song song: (1)~điều khiển PI--ANN thích nghi trên phần cứng thực (ESP32 +
động cơ NF5475), và (2)~huấn luyện bộ điều khiển cân bằng bằng học tăng cường
(RL) trong môi trường mô phỏng NVIDIA Isaac Lab. Cả hai nhánh đều được xây dựng
trên cùng nền tảng cơ khí robot Legged V3 sử dụng cơ cấu 5-bar parallel.


% ==============================================
\section{Mô tả hệ thống robot Legged V3}
% ==============================================

\subsection{Kiến trúc tổng thể}

Robot Legged V3 là robot hai chân có bánh xe (\textit{wheeled bipedal robot}),
trong đó mỗi chân được dẫn động bởi cơ cấu song song 5-bar. Bánh xe được gắn
tại điểm cuối chân (end-effector), cho phép robot vừa di chuyển bằng bánh vừa
thay đổi chiều cao thân thông qua điều khiển góc khớp.

\begin{figure}[H]
    \centering
    % \includegraphics[width=0.80\textwidth]{figure/legged_v3_overview.png}
    \caption{Kiến trúc tổng thể robot Legged V3}
    \label{fig:robot_overview}
\end{figure}

Hệ thống được chia thành ba cấp:
\begin{itemize}
    \item \textbf{Cấp cơ khí:} Khung nhôm, hai cơ cấu 5-bar parallel, trục bánh xe.
    \item \textbf{Cấp dẫn động:} Bốn động cơ DC NF5475 (hai chủ động mỗi chân)
    kết hợp encoder quang học 200 PPR.
    \item \textbf{Cấp điều khiển:} Vi điều khiển ESP32 WROOM-32, mạch driver
    L298N, giao tiếp Serial/UART với máy tính.
\end{itemize}

\subsection{Cơ cấu 5-bar parallel}

Mỗi chân sử dụng cơ cấu 5-bar parallel phẳng với các thông số hình học:
\begin{itemize}
    \item Khoảng cách hai khớp chủ động cố định: $A_1A_2 = 105$ mm.
    \item Chiều dài thanh chủ động: $L_{act} = 100$ mm.
    \item Chiều dài thanh bị động: $L_{pas} = 180$ mm.
\end{itemize}

\begin{figure}[H]
    \centering
    % \includegraphics[width=0.70\textwidth]{figure/five_bar_cad.png}
    \caption{Bản vẽ kỹ thuật cơ cấu 5-bar parallel một chân}
    \label{fig:five_bar_cad}
\end{figure}

Ưu điểm của cơ cấu này so với chân nối tiếp: (i)~độ cứng vững cao hơn do tải
được phân phối trên hai thanh bị động, (ii)~khối lượng phần chuyển động nhỏ vì
động cơ đặt tại gốc cố định, (iii)~giảm số bậc tự do cần điều khiển.

\subsection{Danh sách linh kiện}

\begin{table}[H]
    \centering
    \caption{Danh sách linh kiện hệ thống robot Legged V3}
    \label{tab:bom}
    \begin{tabular}{|c|p{5cm}|c|p{4.5cm}|}
    \hline
    \textbf{STT} & \textbf{Linh kiện} & \textbf{SL} & \textbf{Thông số chính} \\
    \hline
    1 & Động cơ DC NF5475 & 4 & 12--24V, 33W, encoder 200 PPR \\
    \hline
    2 & Vi điều khiển ESP32 WROOM-32 & 1 & 240 MHz, Wi-Fi, BT \\
    \hline
    3 & Mạch driver L298N & 2 & 2A/kênh, 2 chiều \\
    \hline
    4 & Khung nhôm định hình & -- & Cắt CNC theo bản vẽ \\
    \hline
    5 & Nguồn 12V / 5A & 1 & Cấp cho động cơ \\
    \hline
    6 & Module hạ áp LM2596 & 1 & 12V $\to$ 5V cho ESP32 \\
    \hline
    \end{tabular}
\end{table}


% ==============================================
\section{Thiết kế thực nghiệm thu thập dữ liệu}
% ==============================================

\subsection{Quy trình thu thập}

Dữ liệu đặc tính động cơ được thu thập theo giao thức kích thích bậc thang
(step input). ESP32 lần lượt cấp 5 mức PWM: 20\%, 40\%, 60\%, 80\%, 100\%
(tương ứng 51--255 đơn vị 8-bit), mỗi mức duy trì \textbf{10 giây}.

Vận tốc tức thời được tính sau mỗi $\Delta t = 0{,}01$ s:
\begin{equation}
    v_{cur}(k) = \frac{n_{xung}(k)}{200} \cdot \frac{60}{\Delta t}
    \quad [\text{RPM}]
    \label{eq:vel_calc}
\end{equation}
trong đó $n_{xung}(k)$ là số xung encoder tích lũy trong khoảng $\Delta t$.

Dữ liệu $(t,\; v_{cur},\; \text{PWM})$ được xuất qua cổng Serial (115\,200
baud) và lưu thành file CSV trên máy tính. Kết quả thu được 5 file CSV,
tổng khoảng 5\,000 mẫu dữ liệu thô.


% ==============================================
\section{Nhận dạng hàm truyền bằng MATLAB}
% ==============================================

\subsection{Chuẩn bị dữ liệu}

Dữ liệu 5 file CSV được ghép thành một \texttt{iddata} object trong MATLAB với
đầu vào $u$ là mức PWM (chuẩn hóa 0--1) và đầu ra $y$ là vận tốc (RPM),
chu kỳ lấy mẫu $T_s = 0{,}01$ s.

\subsection{Lựa chọn bậc mô hình và nhận dạng}

Hàm \texttt{tfest(data, np, nz)} được thử với $n_p \in \{2,3,4,5\}$ cực và
$n_z \in \{n_p{-}2,\, n_p{-}1\}$ không. Mô hình được chọn theo tiêu chí:
\begin{itemize}
    \item \textbf{Fit-to-data} cao nhất trên tập kiểm chứng độc lập.
    \item \textbf{FPE} và \textbf{MSE} nhỏ.
    \item Không có cực không ổn định (phần thực dương).
\end{itemize}

\subsection{Kiểm chứng mô hình}

Mô hình được kiểm chứng bằng hàm \texttt{compare} trong MATLAB System
Identification Toolbox, so sánh đầu ra mô phỏng với dữ liệu thực nghiệm
trên tập dữ liệu không tham gia nhận dạng.


% ==============================================
\section{Tối ưu tham số PI bằng giải thuật di truyền}
% ==============================================

\subsection{Phát biểu bài toán tối ưu}

Cho hàm truyền $G(s)$ và setpoint $v_{req}$, bài toán là tìm $(K_p, K_i)$ để
cực tiểu hóa hàm thích nghi:
\begin{equation}
    f(K_p, K_i) = w_1 \cdot OS + w_2 \cdot T_s + w_3 \cdot |e_{ss}|
    \label{eq:fitness}
\end{equation}
trong đó $OS$~(\%) là độ vọt lố, $T_s$~(s) là thời gian xác lập, $e_{ss}$
là sai số xác lập, $w_1,w_2,w_3$ là trọng số thực nghiệm.

\subsection{Cấu hình giải thuật di truyền}

\begin{itemize}
    \item \textbf{Chromosome:} $[K_p,\; K_i]$ với $K_p \in [0,\,2]$,
    $K_i \in [0,\,20]$.
    \item \textbf{Quần thể:} 50 cá thể; \textbf{Số thế hệ:} 100.
    \item \textbf{Xác suất lai ghép / đột biến:} 0,8 / 0,05.
\end{itemize}

Quy trình GA được lặp lại độc lập cho 9 mức setpoint từ 250 đến 2250 RPM
(bước 250 RPM), thu được 9 cặp $(K_p^*, K_i^*)$ tương ứng.


% ==============================================
\section{Xây dựng tập dữ liệu và huấn luyện mạng ANN}
% ==============================================

\subsection{Đặc trưng và nhãn}

\begin{itemize}
    \item \textbf{Đầu vào:} sai số vận tốc $e = v_{req} - v_{cur}$ và
    biến thiên sai số $\Delta e = e(k) - e(k{-}1)$.
    \item \textbf{Đầu ra (nhãn):} cặp $(K_p^*, K_i^*)$ tối ưu từ GA tại
    setpoint tương ứng.
    \item \textbf{Phân chia:} 7/9 bộ làm train (với 15\% validation),
    2/9 bộ làm test độc lập.
\end{itemize}

Đặc trưng đầu vào được chuẩn hóa Z-score (thống kê tính chỉ từ train set):
\begin{equation}
    X_{norm} = \frac{X - \mu_{train}}{\sigma_{train} + 10^{-8}}
    \label{eq:zscore}
\end{equation}

\subsection{Kiến trúc và cấu hình huấn luyện}

Mạng ANN được lập trình bằng Python thuần (không dùng framework) với kiến trúc:
\begin{equation}
    \text{Input}(2) \to \text{Dense}(8){+}\text{L2}
    \to \text{ReLU} \to \text{Dropout}(0{,}1)
    \to \text{Dense}(16){+}\text{L2}
    \to \text{ReLU} \to \text{Dropout}(0{,}1)
    \to \text{Dense}(2)
    \label{eq:ann_arch}
\end{equation}

Cấu hình huấn luyện: optimizer Adam ($\alpha = 2\times10^{-3}$, decay
$= 3\times10^{-4}$), hàm loss MSE, 1\,000 epochs, dừng sớm nếu validation
loss không cải thiện sau 50 epochs liên tiếp.


% ==============================================
\section{Triển khai bộ điều khiển PI trên ESP32}
% ==============================================

\subsection{Rời rạc hóa và anti-windup}

Bộ PI được rời rạc hóa bằng Euler tiến ($\Delta t = 0{,}01$ s):
\begin{align}
    e_{vel}(k) &= v_{req}(k) - v_{cur}(k) \\
    I(k) &= \mathrm{clamp}\bigl(I(k{-}1) + e_{vel}(k)\cdot\Delta t,\;
            -1000,\;+1000\bigr) \\
    u(k) &= \mathrm{clamp}\bigl(K_p\cdot e_{vel}(k) + K_i\cdot I(k),\;
            -255,\;+255\bigr)
\end{align}

\subsection{Vòng điều khiển thời gian thực}

Encoder được đọc qua ngắt phần cứng (\texttt{attachInterrupt}, sườn lên kênh A)
để không mất xung. Vòng lặp chính dùng \texttt{millis()} thay cho
\texttt{delay()} để không chặn ISR:
\begin{enumerate}
    \item Kiểm tra $millis() - t_{prev} \geq 10$ ms.
    \item Tính $v_{cur}$ từ bộ đếm xung, reset bộ đếm.
    \item Tính sai số, cập nhật tích phân, tính $u$.
    \item Ghi $u$ ra PWM + chiều quay (L298N).
    \item Xuất log qua Serial; cập nhật $t_{prev}$.
\end{enumerate}


% ==============================================
\section{Huấn luyện điều khiển cân bằng bằng Isaac Lab}
% ==============================================

\subsection{Nền tảng mô phỏng}

Môi trường huấn luyện được xây dựng trên \textbf{NVIDIA Isaac Sim} thông qua
framework \textbf{Isaac Lab}, cho phép chạy song song hàng nghìn môi trường
trên GPU để tăng tốc thu thập dữ liệu. Task được đăng ký với tên
\texttt{Isaac-Legged-V3-Wheel}.

\begin{figure}[H]
    \centering
    % \includegraphics[width=0.85\textwidth]{figure/isaaclab_env.png}
    \caption{Môi trường huấn luyện robot Legged V3 trong NVIDIA Isaac Sim}
    \label{fig:isaaclab_env}
\end{figure}

\subsection{Cấu hình môi trường}

Môi trường sử dụng kiến trúc \texttt{ManagerBasedRLEnv} của Isaac Lab với
4\,096 môi trường song song trên một GPU. Địa hình là mặt phẳng phẳng
(\texttt{terrain\_type = "plane"}). Các cảm biến gồm:
\begin{itemize}
    \item \textbf{IMU} (\texttt{ImuCfg}): gia tốc và vận tốc góc thân robot.
    \item \textbf{Contact Sensor} (\texttt{ContactSensorCfg}): tiếp xúc bánh
    xe với mặt đất.
\end{itemize}

\subsection{Không gian quan sát và hành động}

\textbf{Quan sát} (\textit{observation}) bao gồm:
\begin{itemize}
    \item Lệnh vận tốc đặt $(v_x, v_y, \omega_z)$.
    \item Góc nghiêng và vận tốc góc thân (từ IMU).
    \item Góc và vận tốc góc các khớp chủ động.
    \item Hành động ở bước trước.
\end{itemize}

\textbf{Hành động} (\textit{action}): mô-men xoắn đặt tại các khớp chủ động
của cơ cấu 5-bar và khớp bánh xe.

\subsection{Hàm phần thưởng}

Hàm phần thưởng tổng hợp nhiều thành phần:
\begin{equation}
    r = w_{vel}\,r_{vel} + w_{bal}\,r_{bal}
      + w_{eng}\,r_{eng} + w_{pen}\,r_{pen}
    \label{eq:reward_total}
\end{equation}
trong đó:
\begin{itemize}
    \item $r_{vel}$: thưởng bám vận tốc đặt (theo lệnh $v_x$, $\omega_z$).
    \item $r_{bal}$: thưởng giữ thân thẳng đứng (phạt góc nghiêng lớn).
    \item $r_{eng}$: phạt tiêu thụ năng lượng quá mức.
    \item $r_{pen}$: phạt khi tiếp xúc bất thường hoặc robot ngã.
\end{itemize}

\subsection{Thuật toán và quy trình huấn luyện}

Thuật toán \textbf{PPO} (Proximal Policy Optimization) được sử dụng thông qua
thư viện \texttt{rsl\_rl}. Lệnh huấn luyện:

\begin{lstlisting}[language=bash]
./isaaclab.sh -p scripts/reinforcement_learning/rsl_rl/train.py \
    --task Isaac-Legged-V3-Wheel \
    --num_envs 4096 --headless
\end{lstlisting}

Sau khi huấn luyện, chính sách được kiểm tra bằng lệnh:
\begin{lstlisting}[language=bash]
./isaaclab.sh -p scripts/reinforcement_learning/rsl_rl/play.py \
    --task Isaac-Legged-V3-Wheel \
    --num_envs 4
\end{lstlisting}

\begin{figure}[H]
    \centering
    % \includegraphics[width=0.75\textwidth]{figure/rl_training_curve.png}
    \caption{Đường cong phần thưởng trong quá trình huấn luyện PPO}
    \label{fig:rl_curve}
\end{figure}


% ==============================================
\section{Phân tích động học cơ cấu 5-bar Parallel}
% ==============================================

\subsection{Thiết lập hệ tọa độ}

Phân tích trong mặt phẳng 2D với gốc tọa độ tại $A_1 = (0, 0)$,
$A_2 = (d, 0)$ ($d = 105$ mm), hai góc khớp chủ động $\theta_1$, $\theta_2$
đo từ trục $x$ dương.

\subsection{Động học thuận}

Từ $(\theta_1, \theta_2)$ tính vị trí đầu khâu bị động:
\begin{align}
    B_1 &= \bigl(L_{act}\cos\theta_1,\;\; L_{act}\sin\theta_1\bigr) \\
    B_2 &= \bigl(d + L_{act}\cos\theta_2,\;\; L_{act}\sin\theta_2\bigr)
\end{align}

Điểm cuối $P$ là giao hai đường tròn bán kính $L_{pas}$ tâm $B_1$, $B_2$:
\begin{align}
    (x_P - B_{1x})^2 + (y_P - B_{1y})^2 &= L_{pas}^2 \\
    (x_P - B_{2x})^2 + (y_P - B_{2y})^2 &= L_{pas}^2
\end{align}

\subsection{Động học nghịch}

Từ $P = (x_P, y_P)$ đặt trước, dùng định lý cosin:
\begin{align}
    \theta_1 &= \mathrm{atan2}(B_{1y}, B_{1x})
    - \arccos\!\left(\frac{L_{act}^2 + |A_1P|^2 - L_{pas}^2}
    {2\,L_{act}\,|A_1P|}\right) \\[4pt]
    \theta_2 &= \pi - \mathrm{atan2}(B_{2y}, B_{2x} - d)
    + \arccos\!\left(\frac{L_{act}^2 + |A_2P|^2 - L_{pas}^2}
    {2\,L_{act}\,|A_2P|}\right)
\end{align}

\subsection{Ma trận Jacobian}

Jacobian $J \in \mathbb{R}^{2\times2}$ liên hệ vận tốc khớp với vận tốc
điểm cuối:
\begin{equation}
    \dot{P} = J(\theta)\,\dot{\Theta}, \qquad
    J = \begin{pmatrix}
        \partial x_P/\partial\theta_1 & \partial x_P/\partial\theta_2 \\
        \partial y_P/\partial\theta_1 & \partial y_P/\partial\theta_2
    \end{pmatrix}
    \label{eq:jacobian}
\end{equation}

$J^{-1}$ được dùng để tính vận tốc khớp từ vận tốc điểm cuối mong muốn,
phục vụ bài toán điều khiển quỹ đạo chân trong các giai đoạn sau.
